\documentclass[11pt,a4paper]{article}
\usepackage{turkos}

\begin{document}
\phaseheader{10}{User Mode and System Calls}{Run programs in ring 3, each in its own address space, and let them ask the kernel for services through a single, carefully guarded door}{62}{70}

\ataglance{3 weeks}{5}{Phase 9: locks and wait queues; Phase 6: VMM that works on any page directory; Phase 3: GDT with ring 3 segments}{A TSS, per-process
page directories that share the kernel half, a first jump into ring 3, a system call gate on \code{int 0x80} with a
dispatch table, an ELF loader for programs passed as GRUB modules, a tiny user runtime (\code{crt0}, syscall
wrappers), pointer validation for every user buffer, and faulty programs that get killed while the kernel survives.}

\section{Why this phase}

Everything so far has run in ring 0, where any bug can corrupt anything. An operating system's core promise is
\textbf{protection}: a buggy or malicious program can crash itself, but not the kernel and not other programs. That
promise rests on three pieces of hardware you have already met: privilege rings (Phase 3), paging with the U/S bit
(Phase 6), and the interrupt mechanism as the only way back into ring 0 (Phase 3). This phase connects them.

You will also design the \textbf{system call interface}, the boundary between user programs and the kernel. Every
file operation, every process operation and every byte of output from now on crosses it. Two habits established here
protect you for the rest of the project: every argument from user space is treated as hostile, and every user pointer
is checked before the kernel touches it.

\section{Concepts you need}

\subsection{Limited direct execution}

\OSTEP chapter 6 frames the problem exactly. For speed, user programs should run directly on the CPU. For safety, they
must not be able to do everything. The answer is to run them in \textbf{user mode}, where privileged instructions
(\code{cli}, \code{hlt}, \code{lgdt}, \code{out}, loading CR3) raise \#GP and pages without the U/S bit are
unreachable. When a program needs something only the kernel can do, it executes a \textbf{trap} instruction that
enters the kernel at an address the kernel chose at boot. On x86 that is \code{int 0x80} through an IDT gate whose DPL
is 3, the only gate user code is allowed to invoke. \MOS\,\S1.6 walks through the steps of a \code{read} call, from
the library wrapper down to the kernel and back.

\subsection{What the CPU does on a ring change}

When an interrupt or trap arrives while the CPU is in ring 3, it can't keep using the user's stack: the user controls
that memory, and ESP might point anywhere. So the CPU switches to a kernel stack first, and it finds that stack in the
\textbf{task state segment} (TSS): fields \code{ss0} and \code{esp0}. Then it pushes the user's SS and ESP followed by
the usual EFLAGS, CS and EIP (Figure~\ref{fig:ringchange}); these are the \code{useresp} and \code{ss} fields you left
at the top of \code{struct regs} in Phase 3. \code{iret} reverses the whole thing. Turk-OS uses the TSS only for this;
x86 hardware task switching is ignored. Because each task has its own kernel stack, the scheduler must update
\code{esp0} on every context switch.

\begin{figure}[htbp]
\centering
\begin{tikzpicture}[node distance=0.35cm and 0.5cm, every node/.style={font=\sffamily\scriptsize}]
  \node[tkok, text width=3.0cm] (u) {\textbf{ring 3}: \code{write(1, msg, 5)}\\EAX=4, EBX=1, ECX=msg, EDX=5\\\code{int 0x80}};
  \node[tkbox, text width=3.1cm, right=of u] (g) {IDT entry \code{0x80}: gate DPL 3, so user code may use it};
  \node[tkbox, text width=3.1cm, right=of g] (t) {CPU loads SS:ESP from TSS (\code{ss0}, \code{esp0}) = current task's kernel stack};
  \node[tkbox, text width=3.2cm, below=0.6cm of t] (p) {pushes user SS, ESP, EFLAGS, CS, EIP; jumps to \code{isr128}; CPL = 0};
  \node[tkhi, text width=3.1cm, left=of p] (d) {\code{isr_common} $\rightarrow$ \code{syscall_handler}: \code{table[EAX](EBX, ECX, EDX)}, result in EAX};
  \node[tkok, text width=3.0cm, left=of d] (r) {\code{iret} pops EIP, CS, EFLAGS, ESP, SS: back in ring 3 with EAX = result};
  \draw[tkarrow] (u) -- (g); \draw[tkarrow] (g) -- (t); \draw[tkarrow] (t) -- (p);
  \draw[tkarrow] (p) -- (d); \draw[tkarrow] (d) -- (r);
\end{tikzpicture}
\caption{The round trip of a system call.}
\label{fig:ringchange}
\end{figure}

\subsection{An address space per process}

Each process gets its own page directory. Entries 0--767 (the user half) are private; entries 768--1023 are copied
from the kernel's directory, and because you pre-allocated every kernel page table in Phase 6, those copies stay
correct forever. Switching processes means loading the new directory into CR3. Figure~\ref{fig:uas} shows the layout
of one process. User pages carry the U/S bit in both their PDE and PTE; kernel pages never do, so even though the
kernel is mapped in every address space, user code can't touch it.

\begin{figure}[htbp]
\centering
\begin{tikzpicture}[m/.style={draw=tkNavy!70, minimum width=6.4cm, font=\sffamily\scriptsize, align=center},
  a/.style={font=\ttfamily\scriptsize, anchor=east}]
  \node[m, minimum height=0.8cm, fill=tkRed!10] (k) at (0,0) {kernel half (supervisor only, same in every process)};
  \node[m, minimum height=0.7cm, fill=tkGreen!12, below=0pt of k] (st) {user stack (grows down)};
  \node[m, minimum height=1.4cm, fill=white, below=0pt of st] (gap) {unmapped};
  \node[m, minimum height=0.6cm, fill=tkAmber!12, below=0pt of gap] (hp) {heap, grows up with \code{brk} (Phase 11)};
  \node[m, minimum height=0.9cm, fill=tkBlue!10, below=0pt of hp] (img) {program image from the ELF file:\\\code{.text .rodata .data .bss}};
  \node[m, minimum height=0.6cm, fill=tkGray!15, below=0pt of img] (z) {unmapped: catches \code{NULL} and small offsets};
  \node[a] at ([xshift=-4pt]k.south west) {0xC0000000};
  \node[a] at ([xshift=-4pt]st.south west) {0xBFFFC000};
  \node[a] at ([xshift=-4pt]img.south west) {0x00400000};
  \node[a] at ([xshift=-4pt]z.south west) {0x00000000};
\end{tikzpicture}
\caption{The address space of a Turk-OS user process.}
\label{fig:uas}
\end{figure}

\subsection{Never trust a user pointer}

When a program calls \code{write(1, buf, len)}, \code{buf} is just a number the program chose. It could point into
the kernel (to trick the kernel into printing a secret), be unmapped (to crash the kernel), or wrap around the end of
the address space. Every user pointer must be checked before use: the whole range lies below \code{0xC0000000},
doesn't overflow, and every page in it is present and user-accessible. \code{copy_from_user} and
\code{copy_to_user} do those checks and return \code{-EFAULT} instead of faulting. \MOS\,\S9.7 shows what happens
to kernels that skip this.

\subsection{Loading a program}

A user program is an \textbf{ELF} executable. Its header gives the entry point and a table of \textbf{program
headers}, and each \code{PT_LOAD} header says: take \code{p_filesz} bytes at offset \code{p_offset} in the file, place
them at virtual address \code{p_vaddr}, and zero the rest up to \code{p_memsz} (that tail is \code{.bss}). The loader
allocates and maps pages, copies, zeroes, sets up a user stack, and jumps. \OSC\,\S2.5 covers linkers and loaders.
Until Turk-OS has a file system (Phase 13), programs arrive as \textbf{Multiboot modules}: GRUB (or QEMU's
\code{-initrd}) loads extra files into memory and lists them in the Multiboot information, just as \LOB chapter 7 does.

\section{Reading plan}

\begin{readingplan}
\must & \LOB & Ch.~7 \emph{The Road to User Mode} (GRUB modules, executing a program); Ch.~11 \emph{User Mode}; Ch.~13 \emph{System Calls} & 47--49, 65--68, 71--72 & The x86 steps of Tasks 10.1--10.6. \\
\must & \OSTEP & Ch.~6 \S6.1--6.2 \emph{Basic Technique}, \emph{Restricted Operations} (re-read) & 49--55 & User mode, kernel mode, traps and the trap table. \\
\must & \MOS & \S1.6 \emph{System Calls} (process, file and directory calls; the steps of \code{read}) & 50--62 & What a system call interface looks like, and a model to copy. \\
\should & \OSC & \S2.1--2.5 \emph{Services}, \emph{Interface}, \emph{System Calls}, \emph{Linkers and Loaders} & 55--77 & Parameter passing, API versus system call, loaders. \\
\should & \OSC & \S1.4 \emph{Operating-System Operations} (dual mode, timers) & 21--27 & Why the timer is part of protection. \\
\should & \OSDI & \S1.4 \emph{System Calls} & 26--42 & The POSIX calls a UNIX needs, in one tour. \\
\should & \OSDI & \S2.7 \emph{The System Task in MINIX 3} & 192--204 & How MINIX lets user-space servers request privileged work. \\
\could & \MOS & \S12.2 \emph{Interface Design} & 985--993 & Principles for choosing your system calls. \\
\could & \MOS & \S9.7 \emph{Exploiting Software} & 639--657 & Why every user pointer must be checked. \\
\could & \OSC & \S17.3 \emph{Protection Rings} & 669--671 & Rings in general, and why most OSes use only two. \\
\end{readingplan}

Keep the Intel SDM Volume 3A at hand for the TSS format (the task management chapter) and for the stack switch on a
privilege change (the interrupt chapter), and the ELF specification (\emph{Tool Interface Standard, ELF 1.2}) for the
header formats.

\begin{minixbox}
In MINIX, a POSIX call like \code{read} is not a kernel call at all. The library turns it into a message, and the
kernel's only job is to deliver that message to a user-space server (the file system or the process manager) through
\code{sys_call} in \texttt{kernel/proc.c (07400)}. When a server needs something privileged, such as copying data
between address spaces or writing an I/O port, it sends a \emph{kernel call} to the \textbf{system task},
\texttt{kernel/system.c (09700)}, described in \OSDI\,\S2.7. So MINIX has two layers of calls where Turk-OS has one:
your \code{int 0x80} goes straight to a kernel function.
\end{minixbox}

\section{Build it}

\task{Install a TSS}
Define the 104-byte TSS structure, put it in GDT entry 5 (selector \code{0x28}) with access byte \code{0x89} (present,
ring 0, 32-bit available TSS) and byte granularity, and load it with \code{ltr}. Only \code{ss0} and \code{esp0}
matter. Setting \code{iomap_base} to the structure's size means ``no I/O permission bitmap'', so any port access from
ring 3 faults.

\begin{lst}{c}{kernel/arch/i386/tss.c}
struct tss {
    uint32_t prev_tss;
    uint32_t esp0;              /* kernel stack pointer loaded on entry from ring 3 */
    uint32_t ss0;               /* kernel stack segment: 0x10                       */
    uint32_t esp1, ss1, esp2, ss2;
    uint32_t cr3, eip, eflags, eax, ecx, edx, ebx, esp, ebp, esi, edi;
    uint32_t es, cs, ss, ds, fs, gs, ldt;
    uint16_t trap, iomap_base;
} __attribute__((packed));
_Static_assert(sizeof(struct tss) == 104, "TSS layout");

static struct tss tss;

void tss_init(void)
{
    tss.ss0 = 0x10;
    tss.iomap_base = sizeof(tss);
    gdt_set(5, (uint32_t)&tss, sizeof(tss) - 1, 0x89, 0x0);
    __asm__ volatile ("ltr %0" : : "r"((uint16_t)0x28));
}

void tss_set_kernel_stack(uint32_t esp0) { tss.esp0 = esp0; }
\end{lst}
\verify{\code{info registers} in the QEMU monitor shows \code{TR=0028} with your TSS's base address.}

\task{Give every process an address space}
Add a page directory pointer to \code{struct task} (kernel threads use the kernel's directory). Write
\code{vmm_create_address_space()}: take a page with \code{kalloc_page}, leave entries 0--767 empty, copy entries
768--1023 from the kernel directory. In \code{schedule()}, just before \code{switch_context}, load the next task's
directory into CR3 if it differs from the current one, and call \code{tss_set_kernel_stack} with the top of the next
task's kernel stack.
\verify{Two kernel threads with different directories both run; \code{info mem} shows the kernel mappings in each.}

\task{The first jump into ring 3}
Take a baby step before loading real programs. Create an address space, map one user page at \code{0x00400000}
containing just the two bytes \code{EB FE} (\code{jmp $}, an infinite loop), map one user stack page, and enter ring 3
by building the frame \code{iret} expects: SS, ESP, EFLAGS (with IF set), CS, EIP. The user selectors have RPL 3:
\code{0x1B} for code and \code{0x23} for data.

\begin{lst}{asm}{kernel/arch/i386/usermode.s}
; void enter_user_mode(uint32_t entry, uint32_t user_esp)   -- never returns
global enter_user_mode
enter_user_mode:
    mov  ecx, [esp + 4]     ; user entry point
    mov  edx, [esp + 8]     ; user stack pointer
    mov  ax, 0x23           ; user data selector (GDT index 4, RPL 3)
    mov  ds, ax
    mov  es, ax
    mov  fs, ax
    mov  gs, ax
    push dword 0x23         ; SS
    push edx                ; ESP
    pushfd
    pop  eax
    or   eax, 0x200         ; IF = 1: interrupts stay on in user mode
    push eax                ; EFLAGS
    push dword 0x1B         ; CS (GDT index 3, RPL 3)
    push ecx                ; EIP
    iret                    ; drop to ring 3
\end{lst}
\verify{The monitor stays responsive (the timer preempts the looping user task), and when you stop QEMU while that
task runs, \code{info registers} shows \code{CS=001b} and \code{EIP=00400000}.}

\task{The system call gate}
Add stub 128 to \code{isr.s} and install it with \code{idt_set_gate(0x80, isr128, 0x08, 0xEF)}: a 32-bit
\emph{trap} gate with DPL 3. DPL 3 lets ring 3 use it; the trap type leaves interrupts enabled during system calls,
so a long call can be preempted like any other kernel code (your Phase 9 locks make that safe). Register a handler
that dispatches on EAX. Use the Linux i386 numbers for the calls you implement; then any Linux reference table
matches yours.

\begin{lst}{c}{kernel/syscall/syscall.c}
enum { SYS_exit = 1, SYS_read = 3, SYS_write = 4, SYS_getpid = 20,
       SYS_sched_yield = 158, SYS_nanosleep = 162, NR_SYSCALLS = 256 };

typedef int32_t (*syscall_fn)(uint32_t a, uint32_t b, uint32_t c);

static const syscall_fn syscall_table[NR_SYSCALLS] = {
    [SYS_exit]        = sys_exit,
    [SYS_read]        = sys_read,
    [SYS_write]       = sys_write,
    [SYS_getpid]      = sys_getpid,
    [SYS_sched_yield] = sys_yield,
    [SYS_nanosleep]   = sys_sleep_ms,      /* simplified: argument is milliseconds */
};

static void syscall_handler(struct regs *r)
{
    if (r->eax >= NR_SYSCALLS || !syscall_table[r->eax]) {
        r->eax = (uint32_t)-ENOSYS;
        return;
    }
    r->eax = syscall_table[r->eax](r->ebx, r->ecx, r->edx);   /* return value to user */
}
\end{lst}

\begin{tipbox}[Portability note]
\code{syscall_handler} is generic code, yet it reads x86 registers by name. Dispatch through three one-line inline
helpers in the x86 header instead, such as \code{tf_sysno(r)}, \mbox{\code{tf_arg(r, n)}} and \mbox{\code{tf_set_ret(r, v)}}:
Phase 16 makes exactly these part of the interface every architecture provides. The ports keep these Linux i386 call
numbers and change only the trap: RISC-V uses \code{ecall} with the number in \code{a7} and the arguments in
\code{a0}--\code{a5}, AArch64 uses \texttt{svc~\#0} with the number in \code{x8} and the arguments in
\code{x0}--\code{x5}, and the result comes back in \code{a0} or \code{x0}.
\end{tipbox}

\task{A user runtime}
Create \code{user/} with its own linker script (programs start at \code{0x00400000}), a start file that calls
\code{main} and passes its return value to \code{exit}, and syscall wrappers in inline assembly. Build each program as a
separate ELF with the same cross compiler, linked with \code{-nostdlib} and \code{-T user/user.ld}.

\begin{lst}{asm}{user/crt0.s}
global _start
extern main
extern exit
section .text
_start:
    call main               ; argc and argv arrive in Phase 11
    push eax
    call exit               ; exit(main's return value), never returns
\end{lst}

\begin{lst}{c}{user/lib/syscall.c}
static inline int32_t syscall3(uint32_t n, uint32_t a, uint32_t b, uint32_t c)
{
    int32_t ret;
    __asm__ volatile ("int $0x80" : "=a"(ret) : "a"(n), "b"(a), "c"(b), "d"(c) : "memory");
    return ret;
}

int  write(int fd, const void *buf, unsigned len) { return syscall3(4, fd, (uint32_t)buf, len); }
int  getpid(void)                                  { return syscall3(20, 0, 0, 0); }
void exit(int code)                                { syscall3(1, code, 0, 0); for (;;) { } }
\end{lst}

\task{Load programs as boot modules}
Add \code{module /boot/hello.elf} lines to \code{grub.cfg}. For quick runs, QEMU's Multiboot loader takes the
modules as a comma-separated \code{-initrd} list:

\begin{lst}{sh}{terminal}
qemu-system-i386 -kernel build/kernel.elf -serial stdio \
    -initrd "build/user/hello.elf,build/user/evil.elf"
\end{lst}

In the kernel, walk \code{mods_count} and \code{mods_addr} (through \code{P2V}), reserve the modules' frames in the PMM (you
did that in Phase 5), and keep a table of name, start and end for the loader.

\task{The ELF loader}
Write \code{elf_load(pgdir, image, size, &entry)}. Validate before trusting anything: the magic bytes
\code{7F 45 4C 46}, 32-bit class, little-endian, type \code{ET_EXEC} (2), machine \code{EM_386} (3), and program
headers that lie inside the file. For each \code{PT_LOAD}, check that the segment lies entirely below
\code{0xC0000000} without overflow, allocate and map user pages (writable only if the segment's flags say so), copy
\code{p_filesz} bytes through the direct map, and zero the rest up to \code{p_memsz}. Then map a user stack below
\code{0xC0000000}, create a task that runs \code{enter_user_mode(entry, stack_top)}, and make it ready.

\begin{lst}{c}{include/elf.h}
struct elf32_ehdr {
    uint8_t  e_ident[16];           /* 0x7F 'E' 'L' 'F', class, data, version ... */
    uint16_t e_type, e_machine;
    uint32_t e_version, e_entry, e_phoff, e_shoff, e_flags;
    uint16_t e_ehsize, e_phentsize, e_phnum, e_shentsize, e_shnum, e_shstrndx;
} __attribute__((packed));

struct elf32_phdr {
    uint32_t p_type;                /* 1 = PT_LOAD                  */
    uint32_t p_offset, p_vaddr, p_paddr;
    uint32_t p_filesz, p_memsz;     /* memsz > filesz: zero the tail */
    uint32_t p_flags, p_align;      /* flags: 1 = X, 2 = W, 4 = R    */
} __attribute__((packed));
\end{lst}
\verify{A \code{hello} program prints ``Merhaba from ring 3, pid 3'' through \code{write}, then exits, and its pages go
back to the PMM (free count restored).}

\task{Guard the door}
Write \code{copy_from_user}, \code{copy_to_user} and \code{strncpy_from_user}, and route every system call argument
that is a pointer through them. \code{sys_write} copies the user buffer in chunks into a kernel buffer before printing.
Then make exceptions depend on who caused them: if \code{(r->cs & 3) == 3}, the fault came from user mode, so print
something like ``hello (pid 3): segmentation fault at 0x00000000'' and kill that task; only faults in ring 0 still
panic.

\begin{lst}{c}{kernel/syscall/uaccess.c}
int copy_from_user(void *dst, const void *usrc, size_t n)
{
    uintptr_t start = (uintptr_t)usrc, end = start + n;
    if (end < start || end > KERNEL_VMA)                /* wraps around, or reaches the kernel */
        return -EFAULT;
    for (uintptr_t pg = start & ~0xFFFu; pg < end; pg += PAGE_SIZE)
        if (!vmm_user_page_ok(current->pgdir, pg))      /* present and U/S set in PDE and PTE */
            return -EFAULT;
    memcpy(dst, usrc, n);                               /* current address space is active */
    return 0;
}
\end{lst}

\task{Write evil programs}
A protection mechanism is only as good as its tests. Write user programs that each try one thing they must not do,
load them all at once, and check that each is killed or gets an error while \code{hello} and the kernel carry on.

\begin{center}
\small\renewcommand{\arraystretch}{1.25}
\begin{tabularx}{\linewidth}{@{}l X l@{}}
\toprule
\tkth{Program} & \tkth{What it tries} & \tkth{Expected result} \\
\midrule
\code{divzero} & divide by zero & killed (\#DE from ring 3) \\
\code{nullptr} & read address 0 & killed (page fault, user mode) \\
\code{kpeek} & read \code{*(int *)0xC0100000} & killed (protection violation) \\
\code{cli} & execute \code{cli} or \code{hlt} & killed (\#GP) \\
\code{badptr} & \code{write(1, (void *)0xC0000000, 16)} & \code{write} returns \code{-EFAULT} \\
\code{hugelen} & \code{write(1, buf, 0xFFFFFFFF)} & \code{-EFAULT} (range wraps around) \\
\code{badsys} & \code{int 0x80} with EAX = 9999 & returns \code{-ENOSYS} \\
\bottomrule
\end{tabularx}
\end{center}
\verify{All seven behave as listed in a single boot, and \code{ps} afterwards shows only the kernel's own tasks.}

\section{Testing and debugging}

\begin{pitfalls}
\#GP on the \code{iret} into ring 3 & Selectors without RPL 3, or GDT entries 3--4 with the wrong DPL & Use \code{0x1B} and \code{0x23}; check access bytes \code{0xFA} and \code{0xF2}. \\
Page fault at the user entry point with error code 5 or 7 (user mode, protection violation) & The PDE or the PTE lacks the U/S bit & Both levels need it; print both entries for the faulting address. \\
Triple fault on the first interrupt or syscall from ring 3 & TSS not loaded, \code{esp0} zero, or the TSS descriptor is wrong & Check \code{TR} in \code{info registers}; set \code{esp0} before entering ring 3. \\
User program never gets preempted & EFLAGS pushed for \code{iret} without IF & OR \code{0x200} into EFLAGS in \code{enter_user_mode}. \\
Second process crashes after a context switch & CR3 switched, but \code{esp0} still points at the first task's stack & Update \code{esp0} and CR3 together, on every switch. \\
Kernel crashes inside a system call on a bad pointer & A pointer was used directly instead of through \code{copy_from_user} & Search for every user pointer dereference; there should be none outside \code{uaccess.c}. \\
\end{pitfalls}

\section{Milestone}

\begin{milestonebox}[Protection works]
\begin{checklist}
  \item The TSS is loaded, and \code{esp0} and CR3 are updated on every context switch.
  \item Each process has its own page directory sharing the kernel half.
  \item ELF programs from boot modules run in ring 3 and use \code{write}, \code{getpid}, \code{sleep} and \code{exit}.
  \item All user pointers pass through \code{copy_from_user}/\code{copy_to_user}.
  \item All seven evil programs are contained; the kernel never panics because of user code.
  \item \code{docs/syscalls.md} lists each call: number, arguments, return value, errors.
\end{checklist}
\end{milestonebox}

\begin{questionbox}
\begin{enumerate}
  \item List, in order, everything the CPU does when a ring-3 program executes \code{int 0x80}.
  \item Why does the TSS need \code{esp0}, and why must it change on every context switch?
  \item Why is the system call gate a trap gate with DPL 3, while exception gates are interrupt gates with DPL 0?
  \item How does a system call's return value get back to the user program?
  \item Why must both the PDE and the PTE have the U/S bit set? What happens if only the PTE has it?
  \item What does \code{copy_from_user} check, and what could a program do if it didn't?
  \item The kernel is mapped in every address space. Why can't user code read it?
\end{enumerate}
\end{questionbox}

\section{Going further}

Replace \code{int 0x80} with the faster \code{sysenter}/\code{sysexit} pair and measure the difference with
\code{rdtsc}. Put \code{argc}, \code{argv} and the environment on the initial user stack in the System V i386 layout,
so \code{main(int argc, char **argv)} works (Phase 11 needs this for \code{exec} anyway). Add a read-only page shared by
all processes that holds the current tick count, so programs can read the time without a system call; Linux's vDSO
grew from the same idea.

\nextphase{11}{Process API: fork, exec, wait and exit}
\booklegend
\end{document}
